\documentclass[11pt,a4paper]{article}
\usepackage[utf8]{inputenc}
\usepackage[T1]{fontenc}
\usepackage[spanish,es-minimal,es-noshorthands]{babel}
\usepackage{lmodern}
\usepackage[sfdefault,lining]{FiraSans}
\usepackage[a4paper,top=2.3cm,bottom=2.2cm,left=1.8cm,right=1.8cm,headheight=1.3cm,headsep=0.4cm,footskip=1cm]{geometry}
\usepackage{amsmath,amssymb}
\allowdisplaybreaks
\usepackage{icomma}
\usepackage{graphicx}
\usepackage[table]{xcolor}
\usepackage{tikz}
\usepackage{fontawesome5}
\usepackage{ccicons}
\usepackage{fancyhdr}
\usepackage{lastpage}
\usepackage{tasks}
\usepackage{needspace}
\usepackage[most]{tcolorbox}
\usepackage{array}
\usepackage[hidelinks]{hyperref}
\AtBeginDocument{\hypersetup{pdftitle={Ficha de ejercicios de la unidad 1: Matrices y determinantes. 2º Bach.},pdfsubject={Matemáticas II 2º Bach.},pdfkeywords={matrices, matrices idempotentes, involutivas, nilpotentes, potencias de matrices, determinantes, propiedades de los determinantes, Vandermonde, rango con parámetros, matriz inversa, ecuaciones matriciales, sistemas matriciales}}}

\definecolor{navy}{HTML}{1B2748}
\definecolor{accent}{HTML}{E8603A}
\definecolor{teal}{HTML}{1F9E90}
\definecolor{soft}{HTML}{EEF1F8}
\definecolor{muted}{HTML}{5D6685}

\setlength{\parindent}{0pt}
\setlength{\parskip}{3pt}
\settasks{label=\alph*), label-format=\bfseries\color{navy}, label-width=1.3em, item-indent=1.8em, column-sep=0.6em, after-item-skip=5pt, before-skip=3pt, after-skip=4pt}

% ---------- capçalera i peu ----------
\newcommand{\llicencia}{\href{https://creativecommons.org/licenses/by-nc-sa/4.0/deed.es}{\ccbyncsa}\ \ Licencia Creative Commons Atribución-NoComercial-CompartirIgual 4.0 Internacional (CC BY-NC-SA 4.0)}
\pagestyle{fancy}
\fancyhf{}
\fancyhead[L]{\raisebox{-0.35\height}{\includegraphics[height=0.95cm]{logo-gva.png}}}
\fancyhead[C]{\small\color{muted}Ficha de ejercicios · Unidad 1: Matrices y determinantes · 2\textsuperscript{o} Bach.}
\fancyhead[R]{\raisebox{-0.35\height}{\includegraphics[height=0.95cm]{logo-vermellar.png}}}
\fancyfoot[L]{\scriptsize\color{muted}\llicencia}
\fancyfoot[R]{\small\color{muted}\thepage}
\renewcommand{\headrulewidth}{0.4pt}
\renewcommand{\headrule}{\color{muted!40}\hrule height \headrulewidth}
\fancypagestyle{primera}{\fancyhf{}\renewcommand{\headrulewidth}{0pt}
  \fancyfoot[L]{\scriptsize\color{muted}\llicencia}
  \fancyfoot[R]{\small\color{muted}\thepage}}

% ---------- nivells i exercicis ----------
\newcommand{\estrelles}[1]{{\color{accent}\ifcase#1\or\faStar\faStar[regular]\faStar[regular]\or\faStar\faStar\faStar[regular]\or\faStar\faStar\faStar\fi}}
\newcounter{exer}
\newcommand{\ex}[1]{\par\needspace{4\baselineskip}\medskip\refstepcounter{exer}\noindent{\color{navy}\bfseries\theexer.}\ \ }

\newcommand{\seccio}[2]{\par\needspace{10\baselineskip}\bigskip
  \begin{tcolorbox}[colback=navy,colframe=navy,coltext=white,arc=2mm,boxrule=0pt,left=3mm,right=3mm,top=1.2mm,bottom=1.2mm]
  {\large\bfseries #1\quad #2}
  \end{tcolorbox}\vspace{-2pt}}


% quadre de recordatori (només a la fitxa d'exercicis)
% etiqueta dels exercicis de la PAU
\newcommand{\pau}[1]{{\color{teal}\small\bfseries\faUniversity\ PAU · #1.}\ }
\newcommand{\recorda}[1]{\par\smallskip\begin{tcolorbox}[colback=accent!7,colframe=accent!55,arc=2mm,boxrule=0.5pt,left=2.5mm,right=2.5mm,top=1mm,bottom=1mm]\small{\color{accent}\faLightbulb[regular]}\ \textbf{Recuerda:}\ #1\end{tcolorbox}}
\newcommand{\res}[1]{#1}\newcommand{\rest}[1]{#1}

% =====================================================================
% Capçalera FentMates de les fitxes: un color per a cada curs.
% Els colors són els de la pàgina inicial (LEVELS a assets/js/materials.js);
% el logo ve de _eines/logo/pdf/ (python3 _eines/logo/genera.py).
%
% Ús, al final del preàmbul de la unitat:
%   % =====================================================================
% Capçalera FentMates de les fitxes: un color per a cada curs.
% Els colors són els de la pàgina inicial (LEVELS a assets/js/materials.js);
% el logo ve de _eines/logo/pdf/ (python3 _eines/logo/genera.py).
%
% Ús, al final del preàmbul de la unitat:
%   % =====================================================================
% Capçalera FentMates de les fitxes: un color per a cada curs.
% Els colors són els de la pàgina inicial (LEVELS a assets/js/materials.js);
% el logo ve de _eines/logo/pdf/ (python3 _eines/logo/genera.py).
%
% Ús, al final del preàmbul de la unitat:
%   \input{../comu/fentmates.tex}
%   \fmcurs{1eso}
% i, al començament del document:
%   \capcalera{Matemàtiques}{1r ESO}{Fitxa d'exercicis · Unitat 2}{Divisibilitat}
% Opcional abans de \input: \def\fmneed{14} (espai mínim abans de cada secció).
% =====================================================================
\makeatletter
\def\fm@logodir{../../logo/pdf}
\providecommand{\fmneed}{10}
\def\fm@cursnom{}

% clau del curs, color i color del text damunt del color
\newcommand{\fm@curs}[3]{\@namedef{fm@color@#1}{#2}\@namedef{fm@text@#1}{#3}}
\fm@curs{1eso}{E6B84A}{navy}
\fm@curs{2eso}{3B82F6}{white}
\fm@curs{3eso}{FF7A50}{white}
\fm@curs{4eso-a}{A78BFA}{white}
\fm@curs{4eso-b}{6D28D9}{white}
\fm@curs{1batx-mi}{9B1C1C}{white}
\fm@curs{1batx-ccss}{EF6B6B}{white}
\fm@curs{2batx-mii}{1E6B47}{white}
\fm@curs{2batx-ccss}{2BB3A3}{white}

\newcommand{\fmcurs}[1]{%
  \@ifundefined{fm@color@#1}{\PackageError{fentmates}{Curs desconegut: #1}{}}{}%
  \edef\fm@tmp{\noexpand\definecolor{curs}{HTML}{\@nameuse{fm@color@#1}}%
               \noexpand\colorlet{curstext}{\@nameuse{fm@text@#1}}}\fm@tmp
  \colorlet{cursfosc}{curs!72!black}%
  \def\fm@logo{\fm@logodir/vertical-#1.pdf}%
  \def\fm@marca{\fm@logodir/marca-#1.pdf}}

% ---------- capçalera de la primera pàgina ----------
% #1 matèria (petit), #2 curs, #3 tipus de document i unitat, #4 títol
\newcommand{\capcalera}[4]{%
  \gdef\fm@cursnom{#2}%
  \par\noindent
  \begin{tcolorbox}[colback=curs,colframe=curs,boxrule=0pt,arc=2.2mm,width=3.5cm,height=2.15cm,
      valign=center,left=3mm,right=2mm,top=1mm,bottom=1mm,nobeforeafter,box align=center]
    \color{curstext}{\footnotesize #1}\par\vspace{1pt}{\fontsize{19}{21}\selectfont\bfseries #2\par}%
  \end{tcolorbox}\hspace{5mm}%
  \begin{minipage}[c]{\dimexpr\linewidth-3.5cm-5mm-2.9cm\relax}
    {\color{cursfosc}\small\bfseries\MakeUppercase{#3}}\par\vspace{2pt}
    {\color{navy}\fontsize{20}{23}\selectfont\bfseries #4\par}
  \end{minipage}\hfill
  \begin{minipage}[c]{2.6cm}\raggedleft\includegraphics[width=2.6cm]{\fm@logo}\end{minipage}\par
  \vspace{2pt}}

% ---------- barres de secció del color del curs ----------
\renewcommand{\seccio}[2]{\par\needspace{\fmneed\baselineskip}\bigskip
  \begin{tcolorbox}[colback=curs,colframe=curs,coltext=curstext,arc=2mm,boxrule=0pt,left=3mm,right=3mm,top=1.2mm,bottom=1.2mm]
  {\large\bfseries #1\quad #2}
  \end{tcolorbox}\vspace{-2pt}}

% ---------- estrelles de dificultat del color del curs (a Batxillerat no se'n posen) ----------
\@ifundefined{estrelles}{}{%
  \renewcommand{\estrelles}[1]{{\color{curs}\ifcase#1\or\faStar\faStar[regular]\faStar[regular]\or\faStar\faStar\faStar[regular]\or\faStar\faStar\faStar\fi}}}

% ---------- capçalera i peu de les pàgines ----------
% primera pàgina: logos institucionals al peu
\fancypagestyle{primera}{\fancyhf{}\renewcommand{\headrulewidth}{0pt}
  \fancyfoot[L]{\raisebox{-0.32\height}{\includegraphics[height=0.78cm]{logo-gva.png}}\hspace{5mm}%
    \raisebox{-0.32\height}{\includegraphics[height=0.78cm]{logo-vermellar.png}}}
  \fancyfoot[R]{\scriptsize\color{muted}\href{https://creativecommons.org/licenses/by-nc-sa/4.0/deed.ca}{\ccbyncsa}\ \ CC BY-NC-SA 4.0 · fentmates.org\qquad\small\thepage}}
% la resta: el símbol a l'esquerra i el curs a la dreta, amb el títol al centre com abans
\fancyhead[L]{\raisebox{-0.3\height}{\includegraphics[height=0.62cm]{\fm@marca}}}
\fancyhead[R]{\small\bfseries\color{cursfosc}\fm@cursnom}
\renewcommand{\headrule}{\color{curs!45}\hrule height \headrulewidth}
\makeatother

%   \fmcurs{1eso}
% i, al començament del document:
%   \capcalera{Matemàtiques}{1r ESO}{Fitxa d'exercicis · Unitat 2}{Divisibilitat}
% Opcional abans de \input: \def\fmneed{14} (espai mínim abans de cada secció).
% =====================================================================
\makeatletter
\def\fm@logodir{../../logo/pdf}
\providecommand{\fmneed}{10}
\def\fm@cursnom{}

% clau del curs, color i color del text damunt del color
\newcommand{\fm@curs}[3]{\@namedef{fm@color@#1}{#2}\@namedef{fm@text@#1}{#3}}
\fm@curs{1eso}{E6B84A}{navy}
\fm@curs{2eso}{3B82F6}{white}
\fm@curs{3eso}{FF7A50}{white}
\fm@curs{4eso-a}{A78BFA}{white}
\fm@curs{4eso-b}{6D28D9}{white}
\fm@curs{1batx-mi}{9B1C1C}{white}
\fm@curs{1batx-ccss}{EF6B6B}{white}
\fm@curs{2batx-mii}{1E6B47}{white}
\fm@curs{2batx-ccss}{2BB3A3}{white}

\newcommand{\fmcurs}[1]{%
  \@ifundefined{fm@color@#1}{\PackageError{fentmates}{Curs desconegut: #1}{}}{}%
  \edef\fm@tmp{\noexpand\definecolor{curs}{HTML}{\@nameuse{fm@color@#1}}%
               \noexpand\colorlet{curstext}{\@nameuse{fm@text@#1}}}\fm@tmp
  \colorlet{cursfosc}{curs!72!black}%
  \def\fm@logo{\fm@logodir/vertical-#1.pdf}%
  \def\fm@marca{\fm@logodir/marca-#1.pdf}}

% ---------- capçalera de la primera pàgina ----------
% #1 matèria (petit), #2 curs, #3 tipus de document i unitat, #4 títol
\newcommand{\capcalera}[4]{%
  \gdef\fm@cursnom{#2}%
  \par\noindent
  \begin{tcolorbox}[colback=curs,colframe=curs,boxrule=0pt,arc=2.2mm,width=3.5cm,height=2.15cm,
      valign=center,left=3mm,right=2mm,top=1mm,bottom=1mm,nobeforeafter,box align=center]
    \color{curstext}{\footnotesize #1}\par\vspace{1pt}{\fontsize{19}{21}\selectfont\bfseries #2\par}%
  \end{tcolorbox}\hspace{5mm}%
  \begin{minipage}[c]{\dimexpr\linewidth-3.5cm-5mm-2.9cm\relax}
    {\color{cursfosc}\small\bfseries\MakeUppercase{#3}}\par\vspace{2pt}
    {\color{navy}\fontsize{20}{23}\selectfont\bfseries #4\par}
  \end{minipage}\hfill
  \begin{minipage}[c]{2.6cm}\raggedleft\includegraphics[width=2.6cm]{\fm@logo}\end{minipage}\par
  \vspace{2pt}}

% ---------- barres de secció del color del curs ----------
\renewcommand{\seccio}[2]{\par\needspace{\fmneed\baselineskip}\bigskip
  \begin{tcolorbox}[colback=curs,colframe=curs,coltext=curstext,arc=2mm,boxrule=0pt,left=3mm,right=3mm,top=1.2mm,bottom=1.2mm]
  {\large\bfseries #1\quad #2}
  \end{tcolorbox}\vspace{-2pt}}

% ---------- estrelles de dificultat del color del curs (a Batxillerat no se'n posen) ----------
\@ifundefined{estrelles}{}{%
  \renewcommand{\estrelles}[1]{{\color{curs}\ifcase#1\or\faStar\faStar[regular]\faStar[regular]\or\faStar\faStar\faStar[regular]\or\faStar\faStar\faStar\fi}}}

% ---------- capçalera i peu de les pàgines ----------
% primera pàgina: logos institucionals al peu
\fancypagestyle{primera}{\fancyhf{}\renewcommand{\headrulewidth}{0pt}
  \fancyfoot[L]{\raisebox{-0.32\height}{\includegraphics[height=0.78cm]{logo-gva.png}}\hspace{5mm}%
    \raisebox{-0.32\height}{\includegraphics[height=0.78cm]{logo-vermellar.png}}}
  \fancyfoot[R]{\scriptsize\color{muted}\href{https://creativecommons.org/licenses/by-nc-sa/4.0/deed.ca}{\ccbyncsa}\ \ CC BY-NC-SA 4.0 · fentmates.org\qquad\small\thepage}}
% la resta: el símbol a l'esquerra i el curs a la dreta, amb el títol al centre com abans
\fancyhead[L]{\raisebox{-0.3\height}{\includegraphics[height=0.62cm]{\fm@marca}}}
\fancyhead[R]{\small\bfseries\color{cursfosc}\fm@cursnom}
\renewcommand{\headrule}{\color{curs!45}\hrule height \headrulewidth}
\makeatother

%   \fmcurs{1eso}
% i, al començament del document:
%   \capcalera{Matemàtiques}{1r ESO}{Fitxa d'exercicis · Unitat 2}{Divisibilitat}
% Opcional abans de \input: \def\fmneed{14} (espai mínim abans de cada secció).
% =====================================================================
\makeatletter
\def\fm@logodir{../../logo/pdf}
\providecommand{\fmneed}{10}
\def\fm@cursnom{}

% clau del curs, color i color del text damunt del color
\newcommand{\fm@curs}[3]{\@namedef{fm@color@#1}{#2}\@namedef{fm@text@#1}{#3}}
\fm@curs{1eso}{E6B84A}{navy}
\fm@curs{2eso}{3B82F6}{white}
\fm@curs{3eso}{FF7A50}{white}
\fm@curs{4eso-a}{A78BFA}{white}
\fm@curs{4eso-b}{6D28D9}{white}
\fm@curs{1batx-mi}{9B1C1C}{white}
\fm@curs{1batx-ccss}{EF6B6B}{white}
\fm@curs{2batx-mii}{1E6B47}{white}
\fm@curs{2batx-ccss}{2BB3A3}{white}

\newcommand{\fmcurs}[1]{%
  \@ifundefined{fm@color@#1}{\PackageError{fentmates}{Curs desconegut: #1}{}}{}%
  \edef\fm@tmp{\noexpand\definecolor{curs}{HTML}{\@nameuse{fm@color@#1}}%
               \noexpand\colorlet{curstext}{\@nameuse{fm@text@#1}}}\fm@tmp
  \colorlet{cursfosc}{curs!72!black}%
  \def\fm@logo{\fm@logodir/vertical-#1.pdf}%
  \def\fm@marca{\fm@logodir/marca-#1.pdf}}

% ---------- capçalera de la primera pàgina ----------
% #1 matèria (petit), #2 curs, #3 tipus de document i unitat, #4 títol
\newcommand{\capcalera}[4]{%
  \gdef\fm@cursnom{#2}%
  \par\noindent
  \begin{tcolorbox}[colback=curs,colframe=curs,boxrule=0pt,arc=2.2mm,width=3.5cm,height=2.15cm,
      valign=center,left=3mm,right=2mm,top=1mm,bottom=1mm,nobeforeafter,box align=center]
    \color{curstext}{\footnotesize #1}\par\vspace{1pt}{\fontsize{19}{21}\selectfont\bfseries #2\par}%
  \end{tcolorbox}\hspace{5mm}%
  \begin{minipage}[c]{\dimexpr\linewidth-3.5cm-5mm-2.9cm\relax}
    {\color{cursfosc}\small\bfseries\MakeUppercase{#3}}\par\vspace{2pt}
    {\color{navy}\fontsize{20}{23}\selectfont\bfseries #4\par}
  \end{minipage}\hfill
  \begin{minipage}[c]{2.6cm}\raggedleft\includegraphics[width=2.6cm]{\fm@logo}\end{minipage}\par
  \vspace{2pt}}

% ---------- barres de secció del color del curs ----------
\renewcommand{\seccio}[2]{\par\needspace{\fmneed\baselineskip}\bigskip
  \begin{tcolorbox}[colback=curs,colframe=curs,coltext=curstext,arc=2mm,boxrule=0pt,left=3mm,right=3mm,top=1.2mm,bottom=1.2mm]
  {\large\bfseries #1\quad #2}
  \end{tcolorbox}\vspace{-2pt}}

% ---------- estrelles de dificultat del color del curs (a Batxillerat no se'n posen) ----------
\@ifundefined{estrelles}{}{%
  \renewcommand{\estrelles}[1]{{\color{curs}\ifcase#1\or\faStar\faStar[regular]\faStar[regular]\or\faStar\faStar\faStar[regular]\or\faStar\faStar\faStar\fi}}}

% ---------- capçalera i peu de les pàgines ----------
% primera pàgina: logos institucionals al peu
\fancypagestyle{primera}{\fancyhf{}\renewcommand{\headrulewidth}{0pt}
  \fancyfoot[L]{\raisebox{-0.32\height}{\includegraphics[height=0.78cm]{logo-gva.png}}\hspace{5mm}%
    \raisebox{-0.32\height}{\includegraphics[height=0.78cm]{logo-vermellar.png}}}
  \fancyfoot[R]{\scriptsize\color{muted}\href{https://creativecommons.org/licenses/by-nc-sa/4.0/deed.ca}{\ccbyncsa}\ \ CC BY-NC-SA 4.0 · fentmates.org\qquad\small\thepage}}
% la resta: el símbol a l'esquerra i el curs a la dreta, amb el títol al centre com abans
\fancyhead[L]{\raisebox{-0.3\height}{\includegraphics[height=0.62cm]{\fm@marca}}}
\fancyhead[R]{\small\bfseries\color{cursfosc}\fm@cursnom}
\renewcommand{\headrule}{\color{curs!45}\hrule height \headrulewidth}
\makeatother

\fmcurs{2batx-mii}

\begin{document}
\thispagestyle{primera}

\capcalera{Matemáticas II}{2\textsuperscript{o} Bach.}{Ficha de ejercicios · Unidad 1}{Matrices y determinantes}
\vspace{10pt}
Nombre y apellidos: \rule{6.6cm}{0.4pt}\hfill Grupo: \rule{1.6cm}{0.4pt}\hfill Fecha: \rule{2.2cm}{0.4pt}
\vspace{8pt}

\begin{tcolorbox}[colback=soft,colframe=soft,arc=2mm,left=3mm,right=3mm,top=1.5mm,bottom=1.5mm]
\small Resuelve los ejercicios en el cuaderno mostrando todos los pasos y justificando cada afirmación. Recuerda que el producto de matrices \textbf{no} es conmutativo: en las identidades con matrices y en las ecuaciones matriciales, fíjate bien por qué lado multiplicas. En los ejercicios con parámetros, estudia todos los casos. Los ejercicios marcados con \textcolor{teal}{\faUniversity\ PAU} son de las pruebas de acceso a la universidad de distintas comunidades autónomas (2023--2026).
\end{tcolorbox}

% =====================================================
\seccio{1}{Matrices. Definiciones y tipos}

\ex{1} \label{ex:tipus}Escribe las matrices siguientes y di de qué tipo es cada una:\begin{tasks}(1)
  \task $A=(a_{ij})$, cuadrada de orden 3, con $a_{ij}=i^2-j^2$.
  \task $B=(b_{ij})$, cuadrada de orden 4, con $b_{ij}=1$ si $j=i+1$ \ y \ $b_{ij}=0$ en caso contrario.
  \task $C=(c_{ij})$, cuadrada de orden 3, con $c_{ij}=\min(i,j)$.
\end{tasks}

\ex{2} Sea $A$ una matriz cuadrada cualquiera.\begin{tasks}(1)
  \task Demuestra que $A+A^t$ es simétrica y que $A-A^t$ es antisimétrica.
  \task Deduce que toda matriz cuadrada es la suma de una matriz simétrica y de una antisimétrica, y escribe \ $M=\begin{pmatrix}1 & 4 & -2\\0 & 3 & 5\\6 & -1 & 2\end{pmatrix}$ \ de esta manera.
  \task Demuestra que la única matriz que es a la vez simétrica y antisimétrica es la matriz nula.
\end{tasks}\recorda{$(A+B)^t=A^t+B^t$, \ $(kA)^t=kA^t$ \ y \ $(A^t)^t=A$. \ $A$ es simétrica si $A^t=A$ y antisimétrica si $A^t=-A$.}

\ex{2} \label{ex:idem}Una matriz cuadrada $A$ es \emph{idempotente} si $A^2=A$, \emph{involutiva} si $A^2=I$ y \emph{nilpotente} si alguna potencia es la matriz nula (el menor exponente $k$ con $A^k=O$ se llama \emph{índice}). Clasifica las matrices:\begin{tasks}(2)
  \task $P=\begin{pmatrix}2 & -1\\2 & -1\end{pmatrix}$
  \task $Q=\begin{pmatrix}3 & -4\\2 & -3\end{pmatrix}$
  \task $R=\begin{pmatrix}2 & -4\\1 & -2\end{pmatrix}$
  \task $S=\begin{pmatrix}1 & 1 & 3\\5 & 2 & 6\\-2 & -1 & -3\end{pmatrix}$
\end{tasks}\par ¿Cuál es el índice de nilpotencia de la matriz $B$ del ejercicio \ref{ex:tipus}?

\ex{3} Sea $A$ una matriz idempotente.\begin{tasks}(1)
  \task Demuestra que $I-A$ también es idempotente y que $A(I-A)=O$.
  \task Demuestra que $B=2A-I$ es involutiva.
  \task Compruébalo con la matriz $P$ del ejercicio anterior.
\end{tasks}

\ex{2} \pau{Cataluña, ordinaria 2023} Considera las matrices \ $A=\begin{pmatrix}2 & -3 & -5\\-1 & 4 & 5\\1 & -3 & -4\end{pmatrix}$ \ y \ $B=\begin{pmatrix}2 & 2 & 0\\-1 & -1 & 0\\1 & 2 & 1\end{pmatrix}$.\begin{tasks}(1)
  \task Calcula $A\cdot B$ y $B\cdot A$.
  \task Sean $C$ y $D$ dos matrices cuadradas del mismo orden que cumplen $C\cdot D=C$ y $D\cdot C=D$. Demuestra que $C$ y $D$ son idempotentes.
\end{tasks}


% =====================================================
\seccio{2}{Operaciones con matrices}

\ex{1} Dadas \ $A=\begin{pmatrix}1 & -2 & 0\\3 & 1 & -1\end{pmatrix}$, \ $B=\begin{pmatrix}2 & 1\\0 & -1\\4 & 3\end{pmatrix}$ \ y \ $C=\begin{pmatrix}1 & 2\\-1 & 0\end{pmatrix}$:\begin{tasks}(1)
  \task Calcula $(AB)^t-2C$ y $A\cdot B\cdot C$.
  \task Calcula $B\cdot A$. ¿Tiene sentido comparar $AB$ y $BA$?
  \task Comprueba que $(AB)^t=B^tA^t$ y explica por qué, en general, $A^tB^t$ no es igual a $(AB)^t$ (en este caso, ni siquiera tienen el mismo orden).
\end{tasks}

\ex{2} Sean $A$ y $B$ matrices cuadradas del mismo orden.\begin{tasks}(1)
  \task Demuestra que $(A+B)^2=A^2+2AB+B^2$ si y solo si $AB=BA$.
  \task Demuestra que $(A+B)(A-B)=A^2-B^2$ si y solo si $AB=BA$.
  \task Halla todas las matrices $X$ de orden 2 que conmutan con \ $A=\begin{pmatrix}1 & 2\\0 & 1\end{pmatrix}$, es decir, tales que $AX=XA$.
\end{tasks}

\ex{2} Halla todas las matrices de la forma \ $X=\begin{pmatrix}a & b\\ 0 & c\end{pmatrix}$ \ que cumplen $X^2=I$.

\ex{2} Halla los valores de $m$ para los que \ $A=\begin{pmatrix}1 & m\\ m & 1\end{pmatrix}$ \ cumple $A^2=2A$. Comprueba después que \ $B=\begin{pmatrix}1 & m\\ 0 & -1\end{pmatrix}$ \ cumple $B^2=I$ para cualquier $m$, y calcula $B^{2027}$.

\ex{2} Dada \ $A=\begin{pmatrix}1 & 1 & 1\\0 & 1 & 1\\0 & 0 & 1\end{pmatrix}$, calcula $A^2$ y $A^3$. Conjetura la expresión de $A^n$, comprueba que si es cierta para $n$ también lo es para $n+1$, y calcula $A^{10}$.\recorda{$1+2+\dots+n=\dfrac{n(n+1)}{2}$.}

\ex{2} Dada \ $A=\begin{pmatrix}a & 1\\ 0 & a\end{pmatrix}$, calcula $A^2$, $A^3$ y $A^4$, deduce la expresión de $A^n$ y calcula $A^{10}$ para $a=2$.

\ex{3} \label{ex:J}Sea \ $J=\begin{pmatrix}0 & -1\\1 & 0\end{pmatrix}$.\begin{tasks}(1)
  \task Calcula $J^2$, $J^3$ y $J^4$.
  \task Calcula $J^{2026}$.
  \task Calcula la suma \ $J+J^2+J^3+\dots+J^{2026}$.
\end{tasks}

\ex{3} Sea \ $A=\begin{pmatrix}3 & -4\\1 & -1\end{pmatrix}$ \ y sea $N=A-I$.\begin{tasks}(1)
  \task Comprueba que $N^2=O$.
  \task Demuestra que $A^n=I+nN$ para todo $n\geq1$ (multiplica $I+nN$ por $A=I+N$).
  \task Calcula $A^{50}$.
\end{tasks}

\ex{3} \label{ex:pq}Sea \ $A=\begin{pmatrix}1 & 2\\3 & 4\end{pmatrix}$.\begin{tasks}(1)
  \task Halla los números $p$ y $q$ tales que $A^2=pA+qI$.
  \task Expresa $A^3$ y $A^4$ en la forma $\alpha A+\beta I$, sin calcularlas directamente.
\end{tasks}

\ex{3} La \emph{traza} de una matriz cuadrada, $\operatorname{tr}(A)$, es la suma de los elementos de la diagonal principal.\begin{tasks}(1)
  \task Comprueba que $\operatorname{tr}(AB)=\operatorname{tr}(BA)$ para dos matrices cualesquiera de orden 2.
  \task Deduce que no existe ningún par de matrices de orden 2 tales que $AB-BA=I$.
\end{tasks}

\ex{3} \pau{Madrid, extraordinaria 2025} Dada la matriz \ $A=\begin{pmatrix}1 & 5\\-1 & -3\end{pmatrix}$:\begin{tasks}(1)
  \task Halla las matrices simétricas $B$ que verifican $BA=(A+A^2)B$.
  \task Con $A_1=A$, se consideran las matrices $A_2=A_1^2+A_1$, $A_3=A_2^2+A_2$, $A_4=A_3^2+A_3$, y así sucesivamente. Halla $A_{2025}$.
\end{tasks}


% =====================================================
\seccio{3}{Determinantes}

\ex{1} Calcula el determinante \ $|M|=\begin{vmatrix}1 & 0 & 2 & 0\\3 & -1 & 1 & 2\\0 & 2 & 1 & -1\\2 & 1 & 0 & 3\end{vmatrix}$ \ desarrollando por la primera fila.

\ex{2} Resuelve la ecuación \ $\begin{vmatrix}x & 2 & 3\\ x & x & 3\\ x & x & x\end{vmatrix}=0$.

\ex{3} Considera los determinantes \ $D_2=\begin{vmatrix}m & 1\\ 1 & m\end{vmatrix}$, \ $D_3=\begin{vmatrix}m & 1 & 0\\ 1 & m & 1\\ 0 & 1 & m\end{vmatrix}$ \ y \ $D_4=\begin{vmatrix}m & 1 & 0 & 0\\ 1 & m & 1 & 0\\ 0 & 1 & m & 1\\ 0 & 0 & 1 & m\end{vmatrix}$.\begin{tasks}(1)
  \task Calcula $D_3$ y halla los valores de $m$ que lo anulan.
  \task Desarrollando por la primera fila, comprueba que $D_4=m\cdot D_3-D_2$. Calcula $D_4$ y halla los valores de $m$ que lo anulan.
\end{tasks}

\ex{2} \pau{Navarra, ordinaria 2026} Calcula los valores posibles de $r$ y $t$ para que $|A|+|B|=4$, sabiendo que\[A=\begin{pmatrix}t & t & t\\ t & 3+t & -3+t\\ 2 & -1 & 5\end{pmatrix},\qquad B=\begin{pmatrix}1+r & r & 3r\\ 2 & -1 & 1\\ r & r & 3r\end{pmatrix}\]


% =====================================================
\seccio{4}{Propiedades de los determinantes}

\ex{2} Sabiendo que \ $\begin{vmatrix}a & b & c\\ d & e & f\\ g & h & i\end{vmatrix}=3$, \ calcula, sin desarrollarlos:\begin{tasks}(2)
  \task $\begin{vmatrix}2a & 2b & 2c\\ d-3g & e-3h & f-3i\\ g & h & i\end{vmatrix}$
  \task $\begin{vmatrix}a+d & b+e & c+f\\ d+g & e+h & f+i\\ g+a & h+b & i+c\end{vmatrix}$
  \task $\begin{vmatrix}3c & a & b+2a\\ 3f & d & e+2d\\ 3i & g & h+2g\end{vmatrix}$
  \task $\begin{vmatrix}g & h & i\\ a & b & c\\ d & e & f\end{vmatrix}$
  \task $\begin{vmatrix}a-b & b-c & c-a\\ d-e & e-f & f-d\\ g-h & h-i & i-g\end{vmatrix}$
  \task $|-2A^t|$, donde $A$ es la matriz del determinante dado
\end{tasks}\recorda{Un determinante es lineal respecto a cada fila (o columna): si una fila es una suma, el determinante se descompone en una suma de dos determinantes.}

\ex{2} Sean $A$ y $B$ matrices cuadradas de orden 3 con $|A|=2$ y $|B|=-\tfrac12$, y $C$ una matriz cuadrada de orden 4 con $|C|=-2$. Calcula:\begin{tasks}(4)
  \task $|2A^{-1}B^t|$
  \task $|(3A)^{-1}|$
  \task $|A^2B^{-3}|$
  \task $|A^t(2B)^{-1}A|$
  \task $|BAB^{-1}|$
  \task $|3C|$
  \task $|C^tC^3|$
  \task $|(2C)^{-1}C^2|$
\end{tasks}

\ex{2} Sabiendo que los números 136, 255 y 323 son múltiplos de 17, demuestra, sin calcularlo, que el determinante \ $\begin{vmatrix}1 & 3 & 6\\2 & 5 & 5\\3 & 2 & 3\end{vmatrix}$ \ también es múltiplo de 17. Compruébalo después calculándolo.

\ex{3} \label{ex:lletres}Calcula, haciendo ceros antes de desarrollar, y expresa el resultado factorizado:\begin{tasks}(1)
  \task $\begin{vmatrix}a+b & a & a\\ a & a+b & a\\ a & a & a+b\end{vmatrix}$
  \task $\begin{vmatrix}a & b & b & b\\ b & a & b & b\\ b & b & a & b\\ b & b & b & a\end{vmatrix}$
\end{tasks}

\ex{3} \label{ex:vdm}Demuestra que \ $\begin{vmatrix}1 & 1 & 1\\ a & b & c\\ a^2 & b^2 & c^2\end{vmatrix}=(b-a)(c-a)(c-b)$ \ (\emph{determinante de Vandermonde}) y utilízalo para:\begin{tasks}(1)
  \task calcular \ $\begin{vmatrix}1 & 1 & 1\\ 2 & 3 & 5\\ 4 & 9 & 25\end{vmatrix}$;
  \task resolver la ecuación \ $\begin{vmatrix}1 & 1 & 1\\ 1 & x & 2\\ 1 & x^2 & 4\end{vmatrix}=0$.
\end{tasks}

\ex{3} Razona las cuestiones siguientes:\begin{tasks}(1)
  \task ¿Existe alguna matriz real $A$ de orden 3 tal que $A^2=-I$? ¿Y de orden 2? (Piensa en la matriz $J$ del ejercicio \ref{ex:J}.)
  \task $B$ es una matriz cuadrada de orden 3 tal que $4B^3=B$. ¿Qué valores puede tener $|B|$?
  \task Si $A$ es antisimétrica y de orden \emph{impar}, entonces $|A|=0$.
  \task Si $A$ es de orden 3 y $A\cdot A^t=4I$, ¿cuánto puede valer $|A|$?
\end{tasks}

\ex{2} \pau{Andalucía, ordinaria 2026} Considera la matriz \ $A=\begin{pmatrix}\alpha & \beta & 1\\ 2 & \alpha & \beta\\ 3 & 4 & \alpha\end{pmatrix}$, \ que cumple $|A|=-2$.\begin{tasks}(1)
  \task Calcula \ $\begin{vmatrix}-3 & -4 & -\alpha\\ \alpha-2 & \beta-\alpha & 1-\beta\\ 8 & 4\alpha & 4\beta\end{vmatrix}$.
  \task Calcula \ $\left|A^{-1}\left(A^t\right)^2A\right|$.
\end{tasks}

\ex{2} \pau{Castilla y León, extraordinaria 2026} Si \ $\begin{vmatrix}a & b & c\\ d & e & f\\ g & h & i\end{vmatrix}=-2$, \ calcula el valor de \ $\begin{vmatrix}g-5a & 2d & a\\ 2h-10b & 4e & 2b\\ i-5c & 2f & c\end{vmatrix}$.

\ex{2} \pau{Navarra, extraordinaria 2025} Sean $A$ y $B$ dos matrices cuadradas $3\times3$ tales que $|A|=\tfrac13$ y $|B|=3$. Calcula $|C|$ sabiendo que \ $C=3\cdot\left(A^t\right)^2\cdot(A\cdot B)^{-1}$.

\ex{3} \pau{Asturias, extraordinaria 2025} Una matriz $M$ verifica que $|M|=x$. (Los apartados son independientes.)\begin{tasks}(1)
  \task Supongamos que $M$ es $2\times2$ y que $M^2=(x-1)I$. Calcula todos los valores posibles de $x$.
  \task Supongamos que \ $M=\begin{pmatrix}2 & 1 & 1\\ 0 & -1 & x\\ -1 & -1 & 0\end{pmatrix}$. ¿Existe algún valor de $x$ para el que sea $|M|=x$?
  \task Supongamos que $M$ es $3\times3$, que $x\neq0$ y que $M=xM^2$. Calcula los valores posibles de $x$ y de $|M^{-1}|$.
\end{tasks}

\ex{2} \pau{Murcia, extraordinaria 2024} Una matriz cuadrada $A$ de orden 2 es \emph{ortogonal} si $A\cdot A^t=I$.\begin{tasks}(1)
  \task Estudia si son ortogonales las matrices \ $\begin{pmatrix}\sqrt3/2 & 1/2\\ -1/2 & \sqrt3/2\end{pmatrix}$ \ y \ $\begin{pmatrix}\sqrt3/2 & 1/2\\ -1/2 & -\sqrt3/2\end{pmatrix}$.
  \task Si $A$ es una matriz ortogonal cualquiera de orden 2, calcula razonadamente su determinante.
  \task Justifica que, si $A$ y $B$ son ortogonales de orden 2, el producto $C=A\cdot B$ también lo es.
\end{tasks}

\ex{3} \pau{Navarra, extraordinaria 2026} Dadas las matrices siguientes, calcula $\left|A\cdot B^{-1}\right|$ en función del parámetro $t$. ¿Para qué valores de $t$ no existe $\left|A\cdot B^{-1}\right|$?\[A=\begin{pmatrix}t+1 & t^2 & 2\\ t+1 & t^2+t-2 & 2-t\\ t+1 & t^2 & 4\end{pmatrix},\qquad B=\begin{pmatrix}t & 3t & t^2-1\\ 0 & t & 2-t\\ t & 4t & t^2\end{pmatrix}\]


% =====================================================
\seccio{5}{Rango de una matriz}

\ex{1} Calcula por el método de Gauss el rango de \ $A=\begin{pmatrix}1 & 2 & 0 & -1 & 3\\2 & 3 & 1 & 0 & 1\\0 & 1 & -1 & -2 & 5\\1 & 1 & 2 & 1 & 0\end{pmatrix}$.

\ex{3} Estudia, utilizando menores, el rango de la matriz \ $B=\begin{pmatrix}m & 1 & 1 & 2\\ 1 & m & 1 & 1\\ 1 & 1 & m & 1\end{pmatrix}$ \ según los valores de $m$.\recorda{Que se anule un menor de orden 3 no quiere decir que el rango sea 2: hay que comprobar que se anulan \textbf{todos} los menores de orden 3 que contienen un menor de orden 2 no nulo.}

\ex{3} Estudia el rango de \ $A=\begin{pmatrix}a & 1 & 1 & 1\\ 1 & a & 1 & 1\\ 1 & 1 & a & 1\\ 1 & 1 & 1 & a\end{pmatrix}$ \ según los valores de $a$. (Aprovecha el ejercicio \ref{ex:lletres}.)

\ex{3} Estudia el rango de \ $M=\begin{pmatrix}1 & -1 & 2 & 1\\2 & 1 & a & 3\\1 & 2 & 1 & b\end{pmatrix}$ \ según los valores de $a$ y $b$.

\ex{2} \pau{Castilla-La Mancha, extraordinaria 2026} Estudia el rango de la matriz siguiente en función del parámetro real $a$: \ $A=\begin{pmatrix}a & a-1 & a^2-a\\ a & 1 & a\\ a & 1 & a-1\end{pmatrix}$.

\ex{3} \pau{País Vasco, extraordinaria 2024} Calcula el rango de la matriz $A$ según los valores del parámetro $m$: \ $A=\begin{pmatrix}1 & 3 & 4 & 1\\ m & 2-m & 2 & 1\\ m & -2 & m-2 & 1\end{pmatrix}$.

\ex{2} \pau{Islas Baleares, extraordinaria 2024} Sea $I_3$ la matriz identidad de orden 3 y \ $A=\begin{pmatrix}1 & 0 & 0\\2 & 0 & 0\\-2 & 1 & 1\end{pmatrix}$.\begin{tasks}(1)
  \task Calcula la matriz $B=3A-kI_3$ en función del parámetro real $k$.
  \task Discute el rango de $B$ según el parámetro $k$.
\end{tasks}

\ex{3} \pau{Cantabria, extraordinaria 2025} Considera las matrices \ $A=\begin{pmatrix}-1 & 1\\ 1 & 2\\ 1 & a\end{pmatrix}$ \ y \ $B=\begin{pmatrix}1 & a & 0\\ 3 & 0 & 2\end{pmatrix}$, con $a\in\mathbb{R}$.\begin{tasks}(1)
  \task Estudia el rango de $AB$ en función de los valores de $a$.
  \task Estudia el rango de $BA$ en función de los valores de $a$.
\end{tasks}

\ex{3} \pau{Asturias, ordinaria 2026} Para cada $a\in\mathbb{R}$ se considera la matriz \ $A=\begin{pmatrix}-1 & 2 & -2\\ 0 & 3 & -4\\ a & 2 & -3\end{pmatrix}$.\begin{tasks}(1)
  \task Calcula todos los valores de $a$ que hacen que $A^2=I$.
  \task Si una matriz $B$ de tamaño $3\times3$ verifica $B^2=2I$, ¿qué valores puede tomar $\det(B)$?
  \task Calcula el rango y el determinante de $A$ para cada valor de $a$.
\end{tasks}


% =====================================================
\seccio{6}{Matriz inversa}

\ex{2} Calcula la inversa de las matrices siguientes:\begin{tasks}(2)
  \task $A=\begin{pmatrix}2 & 1 & -1\\0 & 1 & 3\\1 & 0 & 2\end{pmatrix}$
  \task $B=\begin{pmatrix}1 & a & b\\ 0 & 1 & c\\ 0 & 0 & 1\end{pmatrix}$ \ ($a$, $b$, $c$ reales)
\end{tasks}\recorda{$A^{-1}=\dfrac{1}{|A|}\big(\mathrm{Adj}(A)\big)^t$, y solo existe si $|A|\neq0$.}

\ex{2} Dada \ $A=\begin{pmatrix}m & 1 & 0\\ 0 & m & 1\\ 0 & 0 & m\end{pmatrix}$, halla para qué valores de $m$ tiene inversa y calcúlala en función de $m$.

\ex{2} Dada \ $M=\begin{pmatrix}1 & k & 0\\ k & 1 & k\\ 0 & k & 1\end{pmatrix}$:\begin{tasks}(1)
  \task ¿Para qué valores de $k$ tiene inversa?
  \task Calcula $M^{-1}$ para $k=1$.
\end{tasks}

\ex{3} Sea $N$ una matriz cuadrada tal que $N^3=O$.\begin{tasks}(1)
  \task Demuestra que $I-N$ tiene inversa y que $(I-N)^{-1}=I+N+N^2$.
  \task Utilízalo para calcular la inversa de \ $\begin{pmatrix}1 & -2 & -1\\0 & 1 & -3\\0 & 0 & 1\end{pmatrix}$.
\end{tasks}

\ex{3} Sea \ $A=\begin{pmatrix}0 & 1 & 1\\1 & 0 & 1\\1 & 1 & 0\end{pmatrix}$.\begin{tasks}(1)
  \task Comprueba que $A^2=A+2I$.
  \task Deduce que $A$ es invertible y expresa $A^{-1}$ en función de $A$ y de $I$.
  \task Expresa $A^3$ y $A^4$ en la forma $\alpha A+\beta I$.
  \task Haz el mismo razonamiento para calcular la inversa de la matriz del ejercicio \ref{ex:pq}.
\end{tasks}

\ex{3} Demuestra las propiedades siguientes de matrices invertibles:\begin{tasks}(1)
  \task $(A^t)^{-1}=(A^{-1})^t$.
  \task Si $A$ es simétrica e invertible, $A^{-1}$ también es simétrica.
  \task Una matriz es \emph{ortogonal} si $A\cdot A^t=I$. Comprueba que \ $Q=\frac13\begin{pmatrix}1 & 2 & 2\\2 & 1 & -2\\2 & -2 & 1\end{pmatrix}$ \ es ortogonal y escribe $Q^{-1}$ sin hacer ningún cálculo más.
\end{tasks}

\ex{2} \pau{Canarias, extraordinaria 2026} Dada la matriz \ $A=\begin{pmatrix}3 & -k & 0\\ 1 & 2 & k\\ k & 0 & k\end{pmatrix}$, \ con $k\in\mathbb{R}$:\begin{tasks}(1)
  \task Averigua los valores de $k$ para los que $A$ tiene inversa.
  \task Calcula la inversa de $A$ cuando $k=1$.
\end{tasks}

\ex{2} \pau{Aragón, extraordinaria 2026} Sean $A$, $M$ y $N$ tres matrices cuadradas de orden 4 que satisfacen la condición \ $M\cdot A=I+N\cdot A$.\begin{tasks}(1)
  \task Justifica que $A$ es regular (es decir, invertible).
  \task Si \ $M=\begin{pmatrix}1 & 2 & 3 & 4\\1/2 & 1/2 & \sqrt3 & 0\\1/3 & 5 & 0 & -1\\\pi & 0 & 2 & 2\end{pmatrix}$ \ y \ $N=\begin{pmatrix}0 & 3 & 4 & 2\\1/2 & -1/2 & \sqrt3 & 1\\1/3 & 5 & -1 & 0\\\pi & 0 & 2 & 1\end{pmatrix}$, \ calcula $A^{-1}$.
\end{tasks}

\ex{2} \pau{Cataluña, extraordinaria 2023} Sean \ $A=\begin{pmatrix}2 & 1\\3 & 2\end{pmatrix}$, \ $B=\begin{pmatrix}2 & -1\\-3 & 2\end{pmatrix}$ \ y la matriz identidad $I$ de orden 2.\begin{tasks}(1)
  \task Comprueba que $(A-2I)^2=3I$.
  \task Utilizando esta igualdad, halla la inversa de $A$ en función de $A$ y de $I$, y comprueba que coincide con $B$.
  \task Calcula la matriz $X$ que satisface $A\cdot X=B$.
\end{tasks}

\ex{3} \pau{Murcia, extraordinaria 2025} Sea $A$ una matriz cuadrada de orden 3 que cumple $A^2=O$, donde $O$ es la matriz nula de orden 3.\begin{tasks}(1)
  \task Demuestra que $(A+I)^2=2A+I$ y que $(A+I)^3=3A+I$.
  \task Demuestra que la matriz $I-A$ es la inversa de la matriz $I+A$.
  \task Resuelve la ecuación matricial $X+AX=A$, expresando $X$ en función de $A$.
\end{tasks}

\ex{2} \pau{País Vasco, extraordinaria 2025} Sea $\alpha$ un número real y \ $A=\begin{pmatrix}1 & \alpha & 1\\ 0 & 1 & 1\\ 0 & \alpha & 1\end{pmatrix}$.\begin{tasks}(1)
  \task Halla los valores de $\alpha$ para los que existe la matriz inversa de $A$.
  \task En el caso $\alpha=0$, calcula, si es posible, $A^{-1}$ y $A^{2025}$.
\end{tasks}

\ex{2} \pau{Galicia, extraordinaria 2026} Dadas las matrices \ $A=\begin{pmatrix}1 & 0 & 1\\0 & 2 & 0\end{pmatrix}$ \ y \ $B=\begin{pmatrix}1 & 0 & 1\\0 & 1 & 0\\1 & 0 & 1\end{pmatrix}$:\begin{tasks}(1)
  \task Calcula las matrices $X=AA^t$ e $Y=A^tA$. ¿Son invertibles $X$ o $Y$? Razona la respuesta y calcula las inversas cuando sea posible.
  \task Calcula las matrices de la forma $M=\begin{pmatrix}a\\ b\\ c\end{pmatrix}$ tales que $BM=M$.
\end{tasks}

\ex{3} \pau{La Rioja, extraordinaria 2026} Dada la matriz \ $X=\begin{pmatrix}1 & 1\\1 & 2\\1 & 1\\1 & 3\end{pmatrix}$, \ calcula \ $A=I_4-X\left(X^tX\right)^{-1}X^t$ \ y comprueba que es idempotente, es decir, que $A^2=A$. Demuestra después que esto ocurre para cualquier matriz $X$ para la que $X^tX$ sea invertible.


% =====================================================
\seccio{7}{Ecuaciones y sistemas matriciales}

\ex{2} \label{ex:axb}Resuelve la ecuación $AX+B=2X$, siendo \ $A=\begin{pmatrix}3 & 1 & 0\\0 & 3 & 1\\1 & 0 & 2\end{pmatrix}$ \ y \ $B=\begin{pmatrix}1 & 0\\2 & 1\\0 & -1\end{pmatrix}$.\recorda{Para sacar factor común $X$ hay que escribir $2X=2IX$; y hay que poner el factor por el lado correcto: \ $AX-2IX=(A-2I)X$.}

\ex{2} Resuelve la ecuación $XA+X=B$, siendo \ $A=\begin{pmatrix}0 & 1 & 0\\1 & 1 & 1\\0 & 0 & 0\end{pmatrix}$ \ y \ $B=\begin{pmatrix}1 & 2 & 1\\0 & 1 & -1\end{pmatrix}$. ¿Qué orden debe tener $X$?

\ex{3} Resuelve la ecuación $A\,X^{-1}A=B$, siendo \ $A=\begin{pmatrix}1 & 1\\1 & 2\end{pmatrix}$ \ y \ $B=\begin{pmatrix}2 & 1\\1 & 1\end{pmatrix}$.

\ex{3} Resuelve la ecuación $AXB=C$, siendo \ $A=\begin{pmatrix}1 & 1 & 0\\0 & 1 & 1\\1 & 0 & 0\end{pmatrix}$, \ $B=\begin{pmatrix}2 & 1\\1 & 1\end{pmatrix}$ \ y \ $C=\begin{pmatrix}2 & 2\\1 & 2\\2 & 1\end{pmatrix}$. \ (La inversa de $A$ ya la has calculado en el ejercicio \ref{ex:axb}: es la matriz $A-2I$.)

\ex{3} Sea \ $A=\begin{pmatrix}1 & 0\\2 & 1\\1 & 1\end{pmatrix}$.\begin{tasks}(1)
  \task Resuelve la ecuación $AX=\begin{pmatrix}1 & 2\\3 & 5\\2 & 3\end{pmatrix}$.
  \task Resuelve la ecuación $AX=\begin{pmatrix}1 & 2\\3 & 5\\2 & 4\end{pmatrix}$.
\end{tasks}(¡Atención: $A$ no es cuadrada!)

\ex{3} Dadas \ $A=\begin{pmatrix}1 & 0 & 1\\ 0 & m & 0\\ 1 & 0 & m\end{pmatrix}$ \ y \ $B=\begin{pmatrix}-1\\4\\-2\end{pmatrix}$:\begin{tasks}(1)
  \task ¿Para qué valores de $m$ la ecuación $A^2X=B$ tiene una única solución?
  \task Resuélvela para $m=2$.
\end{tasks}

\ex{3} Sea \ $A=\begin{pmatrix}4 & 1\\2 & 3\end{pmatrix}$ \ y sea $X=\begin{pmatrix}x\\ y\end{pmatrix}$.\begin{tasks}(1)
  \task ¿Para qué valores de $\alpha$ la ecuación $AX=\alpha X$ tiene soluciones distintas de $X=\begin{pmatrix}0\\ 0\end{pmatrix}$?
  \task Halla todas las soluciones para cada uno de estos valores.
  \task Sin calcular $A^{100}$, calcula $A^{100}\begin{pmatrix}1\\ 1\end{pmatrix}$.
\end{tasks}

\ex{3} Resuelve el sistema matricial \ $\left\{\begin{aligned}AX+Y&=\begin{pmatrix}3 & 5\\-1 & 3\end{pmatrix}\\ X-Y&=\begin{pmatrix}-1 & 1\\1 & 1\end{pmatrix}\end{aligned}\right.$ \ siendo \ $A=\begin{pmatrix}1 & 2\\0 & 1\end{pmatrix}$.

\ex{2} \pau{Andalucía, extraordinaria 2026} Considera las matrices \ $A=\begin{pmatrix}2 & 2\\3 & 4\end{pmatrix}$ \ y \ $B=\begin{pmatrix}3 & 1\\2 & 2\end{pmatrix}$.\begin{tasks}(1)
  \task Halla razonadamente el determinante de una matriz $X$ que verifica \ $A^3XB^t=5B^2A^{-1}$.
  \task Determina, si existe, una matriz $Y$ que verifica \ $A^3YB^{-1}=A^2$.
\end{tasks}

\ex{2} \pau{Canarias, ordinaria 2026} Dada la matriz \ $A=\begin{pmatrix}1 & 2 & 0\\0 & 1 & 0\\-1 & 1 & 1\end{pmatrix}$, \ halla dos matrices cuadradas de orden 3 tales que la primera más el doble de la segunda dé la matriz $A$, y la primera menos el doble de la segunda dé la inversa de $A$.

\ex{2} \pau{Castilla-La Mancha, ordinaria 2026}\begin{tasks}(1)
  \task Determina los valores de $a$ y $b$ para que las matrices \ $A=\begin{pmatrix}1 & a & 1\\ 1 & 2 & 1\\ 1 & 1 & b\end{pmatrix}$ \ y \ $B=\begin{pmatrix}3 & -1 & -1\\-1 & 1 & 0\\-1 & 0 & 1\end{pmatrix}$ \ cumplan $A\cdot B=I$.
  \task Resuelve razonadamente el sistema \ $\left\{\begin{aligned}X+Y&=\begin{pmatrix}1 & 0\\2 & 3\end{pmatrix}\\ 2X+3Y&=\begin{pmatrix}4 & 1\\7 & 9\end{pmatrix}\end{aligned}\right.$
\end{tasks}

\ex{2} \pau{Extremadura, extraordinaria 2025} Sean las matrices\[A=\begin{pmatrix}0 & 2 & 3\\2 & 3 & 5\\1 & 3 & 1\end{pmatrix},\qquad B=\begin{pmatrix}-1 & 0 & 1\\0 & -1 & 1\\2 & 1 & -3\end{pmatrix}\]\begin{tasks}(1)
  \task Obtén la inversa de la matriz $A^t+I$, donde $I$ es la matriz identidad de orden 3.
  \task Resuelve la ecuación matricial \ $A^tX-I=2B-X$.
\end{tasks}

\ex{2} \pau{Galicia, ordinaria 2026} Considera la igualdad matricial $A(X-I)B=C$, donde $X$ es cuadrada e $I$ es la matriz identidad.\begin{tasks}(1)
  \task Despeja $X$ suponiendo que $A$ y $B$ son cuadradas e invertibles.
  \task ¿Qué dimensiones tienen $X$ y $B$ si $A$ y $C$ son matrices $1\times2$? Halla $X$, simétrica y con los elementos de la diagonal que suman 0, si $A=(1\ \ 1)$, $B=I$ y $C=(0\ \ 0)$.
\end{tasks}

\ex{3} \pau{Madrid, ordinaria 2026} Dada la matriz real \ $A=\begin{pmatrix}1 & \lambda-2 & \lambda\\ 0 & 1 & \lambda-1\\ 1 & 0 & 1\end{pmatrix}$:\begin{tasks}(1)
  \task Discute el rango de $A$ en función del parámetro $\lambda$.
  \task Para el caso $\lambda=2$, resuelve la ecuación matricial $A^2-AX=I$.
\end{tasks}

\ex{2} \pau{Madrid, extraordinaria 2026} Se consideran las matrices \ $A=\begin{pmatrix}1 & 0 & 1\\0 & 1 & 0\\0 & 2 & 1\end{pmatrix}$, \ $B=\begin{pmatrix}2 & -1\\1 & 0\\-1 & -2\end{pmatrix}$ \ y \ $C=\begin{pmatrix}1 & 0 & 0\\-1 & 1 & 0\\0 & -1 & 1\end{pmatrix}$.\begin{tasks}(1)
  \task Resuelve la ecuación matricial $AX+X=B$.
  \task Calcula el determinante de la matriz $(2C)^{100}$ aplicando las propiedades de los determinantes.
  \task Justifica si las matrices $A$ y $AB$ son invertibles.
\end{tasks}

\ex{2} \pau{Navarra, extraordinaria 2024} Sean $A$, $P$ y $Q$ tres matrices cuadradas regulares tales que $Q\cdot A\cdot P=I$, donde $I$ es la matriz identidad de la misma dimensión.\begin{tasks}(1)
  \task Demuestra que \ $A\cdot P\cdot Q\cdot A=Q^{-1}\cdot P^{-1}$.
  \task Calcula la matriz $A$ en el caso \ $P=\begin{pmatrix}-1 & 1\\2 & -1\end{pmatrix}$ \ y \ $Q=\begin{pmatrix}1 & 0\\1 & 2\end{pmatrix}$.
\end{tasks}

\ex{2} \pau{Islas Baleares, ordinaria 2026} Sean $A$ y $B$ dos matrices $3\times3$ tales que $A$ es invertible y \ $AB+A+I=2A$, \ con $I$ la matriz identidad de orden 3.\begin{tasks}(1)
  \task Calcula la inversa de $A$ en función de $I$ y $B$.
  \task Sean \ $B=\begin{pmatrix}0 & 1 & 2\\1 & 1 & 1\\0 & 3 & 1\end{pmatrix}$ \ y \ $b=\begin{pmatrix}0\\1\\0\end{pmatrix}$. Calcula $x$ tal que $Ax=b$, donde $A$ cumple la igualdad del enunciado.
\end{tasks}

\ex{3} \pau{Cantabria, ordinaria 2026} Considera la matriz \ $A=\begin{pmatrix}1 & a & 0\\ 0 & 1 & 0\\ 0 & 0 & a\end{pmatrix}$, \ que depende del parámetro $a\in\mathbb{R}$.\begin{tasks}(1)
  \task Estudia el rango de $A$ en función del valor del parámetro $a$.
  \task Calcula $A^{100}$ en función del parámetro $a$.
  \task Para $a=-1$, calcula, si es posible, la matriz $X$ que cumple \ $AXA^{-1}=\begin{pmatrix}2 & 2 & 2\\0 & 0 & -1\\0 & -2 & -2\end{pmatrix}$.
\end{tasks}

\ex{3} \pau{Castilla y León, extraordinaria 2026} Sean las matrices\[A=\begin{pmatrix}2 & 2 & 0\\0 & 2 & 0\\2 & 0 & 4\end{pmatrix},\qquad B=\begin{pmatrix}-3 & 4 & -6\\-2 & 1 & -2\\-11 & 3 & -8\end{pmatrix}\]Resuelve la ecuación matricial \ $\tfrac14A^2-AX=B$.

\ex{3} \pau{Andalucía, ordinaria 2024} Considera la matriz \ $A=\begin{pmatrix}1 & 1/8 & 1/8\\ 0 & 1 & 0\\ 0 & 0 & 1\end{pmatrix}$.\begin{tasks}(1)
  \task Calcula $A^{2024}$.
  \task Halla la matriz $X$, si es posible, que verifica $A^2XA+I=O$, donde $I$ y $O$ son la matriz identidad y la matriz nula de orden 3.
\end{tasks}


% =====================================================
\seccio{8}{Problemas y cuestiones}

\ex{2} En una liga de cuatro equipos (A, B, C y D), cada equipo juega una vez contra cada uno de los otros y no hay empates. A gana a B y a C; B gana a C y a D; C gana a D, y D gana a A.\begin{tasks}(1)
  \task Escribe la matriz $M$ en la que $m_{ij}=1$ si el equipo $i$ ha ganado al equipo $j$, y $m_{ij}=0$ si no.
  \task Calcula $M^2$. ¿Qué indica el elemento $(i,j)$ de $M^2$?
  \task Para ordenar los equipos, cada equipo obtiene un punto por cada victoria directa y un punto por cada victoria ``indirecta'' (ganar a un equipo que ha ganado a otro). Utiliza $M+M^2$ para hacer la clasificación.
\end{tasks}

\ex{2} Una economía tiene dos sectores, agricultura (1) e industria (2). Para producir 1 € del sector 1 hacen falta 0,20 € del sector 1 y 0,40 € del sector 2; para producir 1 € del sector 2 hacen falta 0,30 € del sector 1 y 0,10 € del sector 2. Si $A=\begin{pmatrix}0{,}2 & 0{,}3\\0{,}4 & 0{,}1\end{pmatrix}$ y $x$ es el vector de producción, la producción que queda para el consumo es $d=x-Ax=(I-A)x$.\begin{tasks}(1)
  \task Si la producción es de 100 y 80 millones de euros, ¿qué parte se destina al consumo?
  \task ¿Qué producción hace falta para que el consumo sea de 56 y 32 millones? Resuélvelo con $(I-A)^{-1}$.
\end{tasks}

\ex{2} Para codificar un mensaje, se sustituye cada letra por su lugar en el alfabeto (A$=1$, B$=2$, \ldots, Z$=26$; espacio $=0$), se agrupan los números de dos en dos en matrices columna y se multiplica cada una por $H=\begin{pmatrix}2 & 1\\3 & 2\end{pmatrix}$. Hemos recibido\[\begin{pmatrix}27\\41\end{pmatrix},\quad\begin{pmatrix}45\\70\end{pmatrix},\quad\begin{pmatrix}38\\57\end{pmatrix}\]Descifra el mensaje. (A B C D E F G H I J K L M N O P Q R S T U V W X Y Z.)

\ex{3} Di si son verdaderas o falsas las afirmaciones siguientes ($A$, $B$, $C$ son matrices cuadradas del mismo orden). Si son verdaderas, demuéstralo; si son falsas, da un contraejemplo.\begin{tasks}(1)
  \task Si $AB=O$, entonces $A=O$ o $B=O$.
  \task Si $A$ es invertible y $AB=AC$, entonces $B=C$.
  \task $|A+B|=|A|+|B|$.
  \task Si $A^2=A$ y $|A|\neq0$, entonces $A=I$.
  \task Si $AB=I$, entonces también $BA=I$.
\end{tasks}

\ex{2} \pau{Aragón, ordinaria 2026} En un ecosistema interactúan tres especies, A, B y C. La población de cada especie en el mes $k$ es $a_k$, $b_k$ y $c_k$, respectivamente, y la evolución de las poblaciones de un mes al siguiente viene dada por\[\begin{pmatrix}a_{k+1}\\ b_{k+1}\\ c_{k+1}\end{pmatrix}=M\begin{pmatrix}a_k\\ b_k\\ c_k\end{pmatrix},\qquad M=\begin{pmatrix}1 & 0 & 0\\1 & 1 & 0\\1 & 0 & 1\end{pmatrix}\]\begin{tasks}(1)
  \task Si en el instante inicial ($k=0$) hay 100 ejemplares de la especie A, 124 de la B y 234 de la C, ¿cuántos ejemplares de cada especie habrá dos meses después?
  \task Razona (sin calcular $M^k$) que se cumple \ $\begin{pmatrix}a_k\\ b_k\\ c_k\end{pmatrix}=M^k\begin{pmatrix}a_0\\ b_0\\ c_0\end{pmatrix}$.
  \task Calcula $M^k$ para $k\in\mathbb{N}$. ¿Qué pasará con el número de ejemplares de cada especie con el paso del tiempo?
\end{tasks}

\vfill
\begin{center}\color{muted}\small\faCheckCircle[regular]\ \ ¡Ánimo! Comprueba tus resultados con el laboratorio de matrices de los apuntes interactivos.\end{center}

\end{document}
